\section{Discussion and Limits}
\label{sec:discussion}

Certified information repair and useful scheduling adaptation are distinct
claims. The diagnostic probes exhibit requirements that no pointwise
scorer of the base representation can satisfy. In the present temporal
and C3 evidence, however, the base quotient is acyclic. Their certified
preference changes support adaptive scoring but do not establish an
information obstruction. Broader workloads are needed to determine where
representation contradictions arise naturally. Moreover, finite quotient
acyclicity neither guarantees realizability in the bounded rule language
nor establishes consistency on unseen decisions.

The repair optimum is relative to a finite proposal catalogue, fixed
certificates, and an additive standalone-work surrogate. It does not
minimize shared compiled cost or complete-schedule regret. Conditional
preferences depend on explicit boundary commitments; a candidate may
avoid that boundary or select a third action. Thus relational consistency
must accompany reachability, escape, and complete-schedule evaluation.
Budget-limited bounds can remain unresolved, and feasibility extends only
to the resource predicates represented by the conflict graph.

The generation comparison contains one independent proposal batch per
condition. Matched requests and candidate counts do not match compute because
pre-delivery checks differ. Results characterize these procedures, not a
model-level guidance advantage; repeated batches and model families are
needed for that conclusion. Dense intersections and unsupported features
also retain substantial evaluation cost. A frozen program's feasibility is
invariant, but its quality--cost tradeoff remains empirical.
